%%
%% This is file `sample-authordraft.tex',
%% generated with the docstrip utility.
%%
%% The original source files were:
%%
%% samples.dtx  (with options: `authordraft')
%% 
%% IMPORTANT NOTICE:
%% 
%% For the copyright see the source file.
%% 
%% Any modified versions of this file must be renamed
%% with new filenames distinct from sample-authordraft.tex.
%% 
%% For distribution of the original source see the terms
%% for copying and modification in the file samples.dtx.
%% 
%% This generated file may be distributed as long as the
%% original source files, as listed above, are part of the
%% same distribution. (The sources need not necessarily be
%% in the same archive or directory.)
%%
%% Commands for TeXCount
%TC:macro \cite [option:text,text]
%TC:macro \citep [option:text,text]
%TC:macro \citet [option:text,text]
%TC:envir table 0 1
%TC:envir table* 0 1
%TC:envir tabular [ignore] word
%TC:envir displaymath 0 word
%TC:envir math 0 word
%TC:envir comment 0 0
%%
%%
%% The first command in your LaTeX source must be the \documentclass command.

%%%%%%%%%%% for the web conference
% \documentclass[sigconf,authordraft]{acmart}

\documentclass{article} % For LaTeX2e
\usepackage{colm2024_conference}

%% NOTE that a single column version may required for 
%% submission and peer review. This can be done by changing
%% the \doucmentclass[...]{acmart} in this template to 
%% \documentclass[manuscript,screen]{acmart}
%% 
%% To ensure 100% compatibility, please check the white list of
%% approved LaTeX packages to be used with the Master Article Template at
%% https://www.acm.org/publications/taps/whitelist-of-latex-packages 
%% before creating your document. The white list page provides 
%% information on how to submit additional LaTeX packages for 
%% review and adoption.
%% Fonts used in the template cannot be substituted; margin 
%% adjustments are not allowed.

%%
%% \BibTeX command to typeset BibTeX logo in the docs
\AtBeginDocument{%
  \providecommand\BibTeX{{%
    \normalfont B\kern-0.5em{\scshape i\kern-0.25em b}\kern-0.8em\TeX}}}

\newcommand{\benyou}[1]{{ \color{blue} [Benyou says ``#1'']}}

\newcommand{\junying}[1]{{ \color{teal} [Junying says ``#1'']}}
\usepackage{multirow}
\usepackage{tabularx}
\usepackage{makecell}
\usepackage{CJK}
% \usepackage[utf8]{inputenc}
%% Rights management information.  This information is sent to you
%% when you complete the rights form.  These commands have SAMPLE
%% values in them; it is your responsibility as an author to replace
%% the commands and values with those provided to you when you
%% complete the rights form.
\setcopyright{acmcopyright}
\copyrightyear{2018}
\acmYear{2018}
\acmDOI{XXXXXXX.XXXXXXX}

%% These commands are for a PROCEEDINGS abstract or paper.
\acmConference[Conference acronym 'XX]{Make sure to enter the correct
  conference title from your rights confirmation emai}{June 03--05,
  2018}{Woodstock, NY}
%
%  Uncomment \acmBooktitle if th title of the proceedings is different
%  from ``Proceedings of ...''!
%
%\acmBooktitle{Woodstock '18: ACM Symposium on Neural Gaze Detection,
%  June 03--05, 2018, Woodstock, NY} 
\acmPrice{15.00}
\acmISBN{978-1-4503-XXXX-X/18/06}


%%
%% Submission ID.
%% Use this when submitting an article to a sponsored event. You'll
%% receive a unique submission ID from the organizers
%% of the event, and this ID should be used as the parameter to this command.
%%\acmSubmissionID{123-A56-BU3}

%%
%% For managing citations, it is recommended to use bibliography
%% files in BibTeX format.
%%
%% You can then either use BibTeX with the ACM-Reference-Format style,
%% or BibLaTeX with the acmnumeric or acmauthoryear sytles, that include
%% support for advanced citation of software artefact from the
%% biblatex-software package, also separately available on CTAN.
%%
%% Look at the sample-*-biblatex.tex files for templates showcasing
%% the biblatex styles.
%%

%%
%% For managing citations, it is recommended to use bibliography
%% files in BibTeX format.
%%
%% You can then either use BibTeX with the ACM-Reference-Format style,
%% or BibLaTeX with the acmnumeric or acmauthoryear sytles, that include
%% support for advanced citation of software artefact from the
%% biblatex-software package, also separately available on CTAN.
%%
%% Look at the sample-*-biblatex.tex files for templates showcasing
%% the biblatex styles.
%%

%%
%% The majority of ACM publications use numbered citations and
%% references.  The command \citestyle{authoryear} switches to the
%% "author year" style.
%%
%% If you are preparing content for an event
%% sponsored by ACM SIGGRAPH, you must use the "author year" style of
%% citations and references.
%% Uncommenting
%% the next command will enable that style.
%%\citestyle{acmauthoryear}

%%
%% end of the preamble, start of the body of the document source.

\begin{document}

%%
%% The "title" command has an optional parameter,
%% allowing the author to define a "short title" to be used in page headers.
\title{HuatuoGPT-II, One-stage Training for Medical Adaption of LLMs}

%%
%% The "author" command and its associated commands are used to define
%% the authors and their affiliations.
%% Of note is the shared affiliation of the first two authors, and the
%% "authornote" and "authornotemark" commands
%% used to denote shared contribution to the research.
\author{Ben Trovato}
\authornote{Both authors contributed equally to this research.}
\email{trovato@corporation.com}
\orcid{1234-5678-9012}
\author{G.K.M. Tobin}
\authornotemark[1]
\email{webmaster@marysville-ohio.com}
\affiliation{%
  \institution{Institute for Clarity in Documentation}
  \streetaddress{P.O. Box 1212}
  \city{Dublin}
  \state{Ohio}
  \country{USA}
  \postcode{43017-6221}
}

% \author{Lars Th{\o}rv{\"a}ld}
% \affiliation{%
%   \institution{The Th{\o}rv{\"a}ld Group}
%   \streetaddress{1 Th{\o}rv{\"a}ld Circle}
%   \city{Hekla}
%   \country{Iceland}}
% \email{larst@affiliation.org}

% \author{Valerie B\'eranger}
% \affiliation{%
%   \institution{Inria Paris-Rocquencourt}
%   \city{Rocquencourt}
%   \country{France}
% }

% \author{Aparna Patel}
% \affiliation{%
%  \institution{Rajiv Gandhi University}
%  \streetaddress{Rono-Hills}
%  \city{Doimukh}
%  \state{Arunachal Pradesh}
%  \country{India}}

% \author{Huifen Chan}
% \affiliation{%
%   \institution{Tsinghua University}
%   \streetaddress{30 Shuangqing Rd}
%   \city{Haidian Qu}
%   \state{Beijing Shi}
%   \country{China}}

% \author{Charles Palmer}
% \affiliation{%
%   \institution{Palmer Research Laboratories}
%   \streetaddress{8600 Datapoint Drive}
%   \city{San Antonio}
%   \state{Texas}
%   \country{USA}
%   \postcode{78229}}
% \email{cpalmer@prl.com}

% \author{John Smith}
% \affiliation{%
%   \institution{The Th{\o}rv{\"a}ld Group}
%   \streetaddress{1 Th{\o}rv{\"a}ld Circle}
%   \city{Hekla}
%   \country{Iceland}}
% \email{jsmith@affiliation.org}

% \author{Julius P. Kumquat}
% \affiliation{%
%   \institution{The Kumquat Consortium}
%   \city{New York}
%   \country{USA}}
% \email{jpkumquat@consortium.net}

%%
%% By default, the full list of authors will be used in the page
%% headers. Often, this list is too long, and will overlap
%% other information printed in the page headers. This command allows
%% the author to define a more concise list
%% of authors' names for this purpose.
\renewcommand{\shortauthors}{Trovato and Tobin, et al.}

%%
%% The abstract is a short summary of the work to be presented in the
%% article.
\begin{abstract}
There exists a growing attention to adapt Large-scale language models like LLama2 and the Baichuan to 
vertical domain, \textit{a.k.a} `domain adaption'.  Typically, it involves a two-step process including \textit{further pre-training} and \textit{supervised fine-tuning}, which introduces challenges such as inconsistent data distribution and the need for additional hyper-parameters. To address these, we propose a simplified, unified domain adaptation protocol, merging both stages into one. Technically, we transform varied data in different language and genres into straightforward input-output pairs. The challenge is the heterogeneity of training data in language and genres. We employ a 'back-translation' method to leverage diverse data genres and facilitate the use of high-quality multilingual data. 
Experiments says something... 
\end{abstract}

%%
%% The code below is generated by the tool at http://dl.acm.org/ccs.cfm.
%% Please copy and paste the code instead of the example below.
%%
\begin{CCSXML}
<ccs2012>
 <concept>
  <concept_id>00000000.0000000.0000000</concept_id>
  <concept_desc>Do Not Use This Code, Generate the Correct Terms for Your Paper</concept_desc>
  <concept_significance>500</concept_significance>
 </concept>
 <concept>
  <concept_id>00000000.00000000.00000000</concept_id>
  <concept_desc>Do Not Use This Code, Generate the Correct Terms for Your Paper</concept_desc>
  <concept_significance>300</concept_significance>
 </concept>
 <concept>
  <concept_id>00000000.00000000.00000000</concept_id>
  <concept_desc>Do Not Use This Code, Generate the Correct Terms for Your Paper</concept_desc>
  <concept_significance>100</concept_significance>
 </concept>
 <concept>
  <concept_id>00000000.00000000.00000000</concept_id>
  <concept_desc>Do Not Use This Code, Generate the Correct Terms for Your Paper</concept_desc>
  <concept_significance>100</concept_significance>
 </concept>
</ccs2012>
\end{CCSXML}

\ccsdesc[500]{Do Not Use This Code~Generate the Correct Terms for Your Paper}
\ccsdesc[300]{Do Not Use This Code~Generate the Correct Terms for Your Paper}
\ccsdesc{Do Not Use This Code~Generate the Correct Terms for Your Paper}
\ccsdesc[100]{Do Not Use This Code~Generate the Correct Terms for Your Paper}

%%
%% Keywords. The author(s) should pick words that accurately describe
%% the work being presented. Separate the keywords with commas.
\keywords{Do, Not, Us, This, Code, Put, the, Correct, Terms, for,
  Your, Paper}

%% A "teaser" image appears between the author and affiliation
%% information and the body of the document, and typically spans the
%% page.
% \begin{teaserfigure}
%   \includegraphics[width=\textwidth]{sampleteaser}
%   \caption{Seattle Mariners at Spring Training, 2010.}
%   \Description{Enjoying the baseball game from the third-base
%   seats. Ichiro Suzuki preparing to bat.}
%   \label{fig:teaser}
% \end{teaserfigure}

\received{20 February 2007}
\received[revised]{12 March 2009}
\received[accepted]{5 June 2009}

%%
%% This command processes the author and affiliation and title
%% information and builds the first part of the formatted document.
\maketitle


\section{Introduction}

Currently, general large-scale models, such as LLama2 and the Baichuan series, are developing particularly rapidly. However, in some vertical domains, some researchers  attempt to develop specialized models for vertical tasks. Vertical models has the potential to yield results comparable to those of larger models by utilizing a medium-sized model, especially in vertical domains, through the exclusion of certain general knowledge. For instance, knowledge in finance may not be applicable in the medical field and can be omitted in a moderately-sized medical LLM, thereby freeing up more capacity for memorizing medical knowledge. Furthermore, this allows for greater flexibility in implementing vertical customization without concern for any negative impact on other domains.  In the future, these specialized models can be incrementally integrated into larger, general-purpose language models, similar to the divide-and-conquer strategy found in computer science.

Adaptation of general-purpose large-scale language models in vertical domains typically involves two steps:\textbf{ continue pre-training} and \textbf{supervised fine-tuning}. It is posited  by \cite{zhou2023lima} that  a language model predominantly gains its knowledge and capabilities during the pre-training phase, and the fine-tuning process fine-tunes its aptitude in selecting the suitable subdistribution of formats for user interactions.
In the  adaptation of such models for specialized fields like medicine, \textit{further pre-training}  aims to inject specialized knowledge, such as medical expertise, while the latter seeks to activate the ability to follow medical instruction. 
% Recent research has also indicated that such domain adaption using further training~\cite{cheng2023adapting} leads to a drastic drop despite still benefiting fine-tuning evaluation and knowledge probing tests. This inspires us to design a simple protocol for domain adaption. 
However, this two-stage training for vertical adaption presents several issues, including:
\junying{In my options, the issues need to be emphasized and be more plausible:

1.Capacity Degradation: Research shows pre-trained models lose inherent and generalization capabilities during distribution transfer. Thus, most researchers prefer base pre-trained models over chat-adapted versions for fine-tuning. The latter undergoes an initial adaptation, causing the degradation of the original pre-trained model. However, "continue pre-training" (an unsupervised learning) and domain-specific fine-tuning (a supervised learning) have different training objectives and data, and can be seen as two separate fine-tunings or distribution transfers. Obviously, this two-stage adaptation could cause an extra inherent degradation, compared to a single-stage transfer.
% This is why most researchers prefer using base pre-trained models for fine-tuning rather than their chat version. The chat-adapted versions have undergone an initial round of fine-tuning, and further tuning can degrade the model's capabilities. Similarly, "continue pre-training" can be viewed as fine-tuning on specific data, since it is different from the model's original general pre-training. However, "continue pre-training" (an unsupervised learning) and domain-specific fine-tuning (a supervised learning) have different training objectives and data, and can be seen as undergoing two distinct fine-tunings. Thus, this two-stage adaptation can compromise the model's performance compared to traditional one-step fine-tuning.


2. Cumbersome and Inefficient. In a dual-stage training process, continued pre-training aims to inject domain-specific knowledge, with the aspiration that this knowledge can amplify subsequent fine-tuning. However, the stark disparities between the two stages render this knowledge-injection inefficient. The pre-training data, characterized by its lengthy content and issues like punctuation errors, ethical dilemmas, and stylistic chaos, greatly contrasts with the fine-tuning data, which is more concise, formal, and standardized. This difference often results in the pre-training data being sidelined during adaptation, viewed as out-of-distribution. Moreover, continued pre-training introduces additional training stage that request extra training architecture and parameter tuning, complicating domain adaptation further.

3. 

% 2.Significant Discrepancy between Continued Pre-training and Domain Fine-tuning: The essence of continued pre-training is to impart domain-specific knowledge into the model, facilitating the transition to the domain-specific fine-tuning phase. For effective knowledge assimilation, it's vital that the continued pre-training aligns with the domain fine-tuning process(references or examples?). However, there's a clear disparity between these two phases. 1) Data Discrepancy. Primarily, the data in each phase differs. Pre-training data is usually of lower quality than fine-tuning data, leading to issues like punctuation mistakes, ethical concerns, stylistic chaos.  2) Discrepancy in Training Modalities.The training methods for both stages also diverge. While pre-training uses unsupervised autoregression, fine-tuning often adopts a question-masking format. Factors like learning rates, batch sizes, architectures, and data sampling strategies differ considerably between the two, increasing the complexity of the adaptation process.

3. Data Scarcer. Building on the general pre-trained models, continued pre-training is crucial for enhancing domain expertise. This process inherently requires more extensive and high-quality domain-specific content, compared to the original corpus of the foundational model. However, compared to general data,  domain-specific data is limited, with truly high-quality data being even harder to come by. For instance, only 3\% of the C4 dataset relates to medicine, which dwindles to mere 1\% after excluding low-quality content like advertisements. This data scarcity in vertical domains underscores the need for a more effective vertical adaption to embed domain knowledge.


% Additionally, while English dominates in medical research data, other languages lag significantly in such specialized content.

% Although pre-training data is abundant, it doesn't have the knowledge density of fine-tuning data. Historically, the volume of pre-training data made up for its quality issues, but domain-specific data tends to be scarcer.

To tackle the adaptation challenges in the specified domain, this paper proposes one-stage medical adaptation method that consolidates 'continue pre-training' and 'domain fine-tuning' into a single stage. This streamlined approach is designed to minimize unwanted model degradation and to boost the efficiency of integrating pre-trained data. Yet, within this integrated phase, it's imperative to bridge the gap between 'continue pre-training' and 'domain fine-tuning'.
}



\begin{itemize}
    \item  \textbf{Data distribution in two stages are different}.  The distribution of data is inconsistent between pre-training and fine-tuning, necessitating the introduction of a warmup step. More importantly, this complex pipeline protocol affects scalability and flexibility. For instance, modifications in pre-training inevitably require re-execution of fine-tuning.
    \item \textbf{More hyper-parameters needed} The two stages have different hyperparameters, including batch size and learning rate, and even the warmup process.
\end{itemize}

\junying{two-stage training process -> two-stage training for vertical domain adaption. (We don't want to replace the conventional two-stage training process,which is solid, but we want to replace the two-stage traning for vertical domain adaption)}

\junying{\textbf{Two-stage adaptation} is the central topic this paper seeks to address. Within domain adaptation, the model requires more high-quality domain-specific data, especially since the backbone model already possesses certain domain capabilities. The challenge is how to harness the inherent domain competence of the backbone while acquiring more domain knowledge. However, the two-stage training process for vertical domain adaptation can lead to the degradation of the original LLM model (citation needed). We aim to merge these two processes into a single phase. Similar to how PLMs are fine-tuned and immediately adapted, we hope to achieve both domain pre-training and domain fine-tuning in a phase akin to fine-tuning.}



To tackle the two above issues, we propose a simper protocol of domain adaption that unifies the two stage (further pre-training and supervised fine-tuning) in a single stage.   The core recipe is to transform data -—whether structured or unstructured -— into a unified protocol: a straightforward input-output pair, like \cite{raffel2020exploring,yuan2022restructured}. Here, the input represents the given context or instruction that does not execute next-word prediction loss, usually framed as a question, and the output is the expected  response.




The challenge of the proposed method is that the  train data are heterogeneous in two aspects: language and genres. 
There are massive amount of high-quality data in English, which could used to be used in other languages~\cite{zeng2023glm-130b,yang2023baichuan}. Secondly, the data genres varies from  encyclopedia, textbooks, examination to  knowledge bases.

To address the challenge, we use a so-called `back-translation' method to transform data in different genres and languages into a  straightforward input-output pair in a specific language; here we use Chinese and one could extend it to any languages. \benyou{more words to say}



\textbf{Experimental results} \benyou{@junyun, the expected results}


\section{Methodology}
\begin{figure*}[ht]
  \centering
  \resizebox{0.95\textwidth}{!}{
  \includegraphics[width=\textwidth]{./figure/huatuoII.png}
  }
  \caption{\label{fig:method}Schematic of HuatuoGPT-II.  }
\end{figure*}

\benyou{@junyin, please finish it as soon as possible}

\junying{The idea is that there are two sub-chapters here, one is data Integration protocols and the other is One-Step Adaptation Training.}

\subsection{Data Integration Protocol}
\junying{Upcoming addition: A diagram illustrating data numbers flow in pipeline.}
The Data Integration Protocol serves as a method to consolidate data from diverse sources, languages, and formats into a uniform language and format. This not only minimizes the learning curve associated with data differentiation but also enhances the acquisition and transfer of medical knowledge.

Furthermore, this protocol adapts pre-trained text from unsupervised learning objectives into fine-tuning targets. This ensures a seamless integration between the pre-training and fine-tuning phases, effectively narrowing the divide between them. Additionally, it addresses and circumvents the ethical concerns associated with pre-training data during this transformation.



\begin{table}[h]
\centering
\begin{tabular}{|c|c|c|c|}
\hline
\multirow{2}{*}{Language} & \multirow{2}{*}{Type of Data} & \multirow{2}{*}{\# Data} & \multirow{2}{*}{Source} \\
& & & \\
\hline
\multirow{4}{*}{Chinese} & Web Corpus & 640,621 & Wudao, Common Crawl \\
\cline{2-4}
& Books & 1,835,931 & Textbook \\
\cline{2-4}
& Encyclopedia & 411,183 & Online Encyclopedia \footnote{Xiaohe} \\
\cline{2-4}
& Literature & 177,261 & Chinese Literature \\
\hline
\multirow{4}{*}{English} & Web Corpus & 394,490 & C4 \\
\cline{2-4}
& Books & 719,187 & the\_pile\_books3, Textbook \\
\cline{2-4}
& Encyclopedia & 147,059 & Wikipida \\
\cline{2-4}
& Literature & 878,241 & PubMed \\
\hline
\end{tabular}
\caption{\label{tab:pretrained data} Summary of the Pretrained Data}
\end{table}

\paragraph{Data Pipeline for Vertical Field}
To bolster the model's medical proficiency in the medical domain, we devised a domain-specific data pipeline (as detailed in the appendix).  Through this approach, we curated top-tier medical content from both Chinese and English sources, including textbooks, encyclopedias, online materials, and academic papers. Despite processing over 100B of data, we selected just about 5 million of the highest quality records, representing roughly 1.4\% of the initial data. Refer to Table~\ref{tab:pretrained data} for the data distribution and the accompanying figure~\ref{待画} for an in-depth visualization of the pipeline's data processing.


\subsection{One-Step Adaptation Training}



\section{Experiments}

\section{Methodology}
\begin{figure}[ht]
  \centering
  \resizebox{0.45\textwidth}{!}{
  \includegraphics[width=\textwidth]{./figure/loss.pdf}
  }
  \caption{\label{fig:Loss}Training loss of HuatuoGPT-II. The training loss starts with a steady decline from the pre-training phase and later transitions seamlessly into the fine-tuning stage. The whole process was smooth and natural.}
\end{figure}


\subsection{Experimental settings}

Baselines \benyou{there are so many baselines}

\subsection{Evaluation protocol}


\begin{table*}[h]
\centering
\begin{tabular}{lccccccccccccccc}
\toprule
\textbf{Model}              & \textbf{MedQA\_zh}  & \textbf{MedQA\_en} & \textbf{MedMCQA} & \textbf{CMB} & \textbf{CMEexma} & \textbf{MMLU} $\diamondsuit$ & \textbf{CMMLU} $\diamondsuit$ & \textbf{C\_Eval} $\diamondsuit$ & \textbf{TruthfulQA} $\diamondsuit$ \\
\hline
DISC-MedLLM  & 40.22 & 28.67 & 0.86 & 32.47 & 36.62 & 7.30 & 1.77 & 0.00 & 5.45  \\
Qwen-7B-Chat  & 21.8      & 24.7  & 33.0  & 28.3    & 30.9 & 37.0 & 28.6 & 39.5  \\
LLaMA2-7B-Chat  & 26.59      & 30.79  & 37.32  & 24.81    & 25.05 & 45.68 & 31.02 & 00.00 & 40.00  \\
LLaMA2-13B-Chat  & 28.58 & 36.68 & 39.66 & 26.51 & 26.60 & 50.67 & 32.64 & 0.00 & 56.36  \\
ChatGLM3-6B   & 42.79 & 28.75 & 35.91 & 39.81 & 43.21 & 47.21 & 46.97 & 0.00 & 49.09  \\
Baichuan2-7B-Chat  & 45.8      & 38.5  & 56.4  & 38.0    & 52.5 & 53.4 & 45.7 & 65.9  \\
Baichuan2-13B-Chat & 50.0      & 42.5  & 55.4  & 40.5    & 54.0 & 57.5 &\textbf{ 54.3} & 61.0 \\
ChatGPT~(API) & 38.1      & 44.7  & 48.0  & \textbf{50.5}    & 43.5 & \textbf{73.0} & 48.6 & 51.2 \\
HuatuoGPT-I~(7B) & 26.3      & 27.0  & 30.7  & 31.5    & 32.0 & 32.8 & 34.3 & 39.0 \\
HuatuoGPT-II~(7B)   & \textbf{64.2}      & \textbf{47.5}  & \textbf{72.7}  & 49.0    &\textbf{ 61.0 }& 54.0 & 60.0 & \textbf{61.0 }\\
\midrule
\end{tabular}
\caption{Medical benchmark results. Evaluation was done using validation data for MedQA, MedMCQA, and CMB. $\diamondsuit$ signifies extraction of only medical-related multiple-choice questions. $\dag$ indicates a test subset with 200 samples tested.}


\label{table:medical_benchmark}
\end{table*}

\subsection{Medical Benchmark}
    \paragraph{MedQA, MedMCQA, CMB, MMLU, Ceval}
    Detailed description or results regarding the mentioned benchmarks can be inserted here.
    
    \paragraph{Medical Qualification Examination}
    Coverage on recent three years of medical exams, professional pharmacist, USMLE, etc.

\junying{CMExam, Huatuo26M, CMedQA}

\begin{table*}[]\small
\renewcommand{\arraystretch}{1.3} % 0.8 of the original row height
\setlength{\tabcolsep}{3.0pt} % 2.5 points column space
\begin{tabular}{l|cccccc|cccccc|ccccc|c}
\hline
\multicolumn{1}{c|}{\multirow{3}{*}{Model}} & \multicolumn{6}{c|}{\begin{CJK}{UTF8}{gbsn} {中医}\end{CJK}}  & \multicolumn{6}{c|}{\begin{CJK}{UTF8}{gbsn} {临床}\end{CJK}}  & \multicolumn{5}{c|}{\begin{CJK}{UTF8}{gbsn} {药师}\end{CJK}}& \multirow{3}{*}{Avg.} \\ \cline{2-18}
\multicolumn{1}{c|}{}  & \multicolumn{3}{c|}{\begin{CJK}{UTF8}{gbsn} {\begin{CJK}{UTF8}{gbsn} {执业助理医师}\end{CJK}}\end{CJK}} & \multicolumn{3}{c|}{\begin{CJK}{UTF8}{gbsn} {执业医师}\end{CJK}}& \multicolumn{3}{c|}{\begin{CJK}{UTF8}{gbsn} {执业助理医师}\end{CJK}} & \multicolumn{3}{c|}{\begin{CJK}{UTF8}{gbsn} {执业医师}\end{CJK}}& \multicolumn{3}{c|}{\begin{CJK}{UTF8}{gbsn} {执业中药师}\end{CJK}}  & \multicolumn{2}{c|}{\begin{CJK}{UTF8}{gbsn} {执业西药师}\end{CJK}}&\\
\multicolumn{1}{c|}{}  & 2015 & 2016 & \multicolumn{1}{c|}{2017} & 2012 & 2013 & 2016 & 2018 & 2019 & \multicolumn{1}{c|}{2020} & 2018 & 2019 & 2020 & 2017 & 2018 & \multicolumn{1}{c|}{2019} & 2021 & 2022 &\\ \hline
DoctorGLM  & 3.00  & 1.36  & \multicolumn{1}{c|}{3.45}  & 1.83  & 1.83  & 2.09  & 2.56  & 3.28  & \multicolumn{1}{c|}{1.64}  & 1.56  & 2.73  & 1.83  & 1.79  & 2.40  & \multicolumn{1}{c|}{2.24}  & 3.12  & 3.75  & 2.38 \\ 
ChatGLM-Med& 20.67 & 23.18 & \multicolumn{1}{c|}{20.69} & 21.83 & 21.83 & 22.56 & 30.34 & 22.95 & \multicolumn{1}{c|}{29.92} & 18.49 & 20.38 & 24.77 & 16.67 & 20.96 & \multicolumn{1}{c|}{15.25} & 20.83 & 22.50 & 21.99\\
BenTsao & 23.33 & 26.82 & \multicolumn{1}{c|}{17.24} & 19.00 & 19.50 & 22.09 & 18.38 & 24.59 & \multicolumn{1}{c|}{27.46} & 21.60 & 18.91 & 18.12 & 20.24 & 20.36 & \multicolumn{1}{c|}{18.83} & 18.96 & 22.08 & 21.03\\
BianQue-2 & 3.67  & 2.27  & \multicolumn{1}{c|}{3.45}  & 4.17  & 4.50  & 3.95  & 4.70  & 5.33  & \multicolumn{1}{c|}{1.64}  & 3.79  & 5.88  & 3.67  & 7.14  & 4.19  & \multicolumn{1}{c|}{8.97}  & 4.58  & 4.38  & 4.49 \\
ChatMed-Consult  & 20.00 & 17.27 & \multicolumn{1}{c|}{14.66} & 21.33 & 18.17 & 20.00 & 27.78 & 21.31 & \multicolumn{1}{c|}{19.26} & 23.83 & 21.43 & 19.50 & 16.67 & 14.97 & \multicolumn{1}{c|}{17.49} & 22.29 & 24.17 & 20.01\\
MedicalGPT & 25.00 & 24.09 & \multicolumn{1}{c|}{21.55} & 26.33 & 27.00 & 26.98 & 38.89 & 29.10 & \multicolumn{1}{c|}{28.28} & 33.41 & 32.14 & 26.15 & 22.62 & 20.96 & \multicolumn{1}{c|}{20.63} & 24.79 & 24.79 & 26.63\\
DISC-MedLLM & 38.33 & 41.82 & \multicolumn{1}{c|}{ 26.72 } & 36.17 & 38.67 & 35.12 & 49.15 & 36.07 & \multicolumn{1}{c|}{ 41.80 } & 41.43 & 36.55 & 35.78 & 25.00 & 22.16 & \multicolumn{1}{c|}{ 21.97 }   & 25.62 & 25.62 & 33.00 \\
HuatuoGPT & 26.00 & 30.91 & \multicolumn{1}{c|}{32.76} & 31.33 & 26.83 & 30.23 & 32.48 & 33.61 & \multicolumn{1}{c|}{27.46} & 32.74 & 28.57 & 30.05 & 18.45 & 27.54 & \multicolumn{1}{c|}{25.11} & 26.04 & 27.08 & 28.66\\
ChatGLM2-6B& 47.00 & 50.45 & \multicolumn{1}{c|}{43.97} & 44.83 & 48.33 & 48.14 & 62.39 & 49.18 & \multicolumn{1}{c|}{52.87} & 50.11 & 48.53 & 45.64 & 36.90 & 36.53 & \multicolumn{1}{c|}{38.57} & 34.79 & 35.62 & 45.52\\

ChatGLM3-6B & 45.67 & 45.00 & \multicolumn{1}{c|}{ 50.00 } & 46.33 & 45.83 & 46.51 & 50.85 & 48.77 & \multicolumn{1}{c|}{ 43.85 } & 41.65 & 43.91 & 43.58 & 42.26 & 28.14 & \multicolumn{1}{c|}{ 35.43 }   & 38.33 & 35.83 & 43.00 \\

Baichuan-7B-chat & 42.00 & 44.09 & \multicolumn{1}{c|}{43.10} & 39.67 & 43.33 & 41.16 & 54.70 & 48.36 & \multicolumn{1}{c|}{44.26} & 42.76 & 42.65 & 39.45 & 35.71 & 38.32 & \multicolumn{1}{c|}{30.49} & 39.79 & 40.83 & 41.80\\
Baichuan2-13B-Chat & 63.24 & 55.00 & \multicolumn{1}{c|}{58.62} & 58.67 & 59.33 & 58.60 & 64.10 & 63.52 & \multicolumn{1}{c|}{60.66} & 65.27 & 60.08 & 59.86 & 52.38 & 41.31 & \multicolumn{1}{c|}{50.22}   & 50.42 & 46.46 & 54.94\\
LLaMA2-7B-Chat & 24.00 & 22.27 & \multicolumn{1}{c|}{ 23.28 } & 27.17 & 26.00 & 22.33 & 29.06 & 29.92 & \multicolumn{1}{c|}{ 27.87 } & 25.84 & 27.31 & 25.23 & 17.26 & 23.35 & \multicolumn{1}{c|}{ 20.63 }   & 23.96 & 21.04 & 24.00 \\

LLaMA2-13B-Chat & 23.00 & 26.82 & \multicolumn{1}{c|}{ 24.14 } & 27.00 & 27.83 & 25.35 & 29.91 & 29.51 & \multicolumn{1}{c|}{ 27.46 } & 26.50 & 28.99 & 25.69 & 21.43 & 26.95 & \multicolumn{1}{c|}{ 21.08 }   & 24.79 & 22.29 & 25.00 \\

Spark(\begin{CJK}{UTF8}{gbsn} {星火}\end{CJK})  & 43.33 & 38.18 & \multicolumn{1}{c|}{39.66} & 38.33 & 38.17 & 42.56 & 54.70 & 44.67 & \multicolumn{1}{c|}{50.00} & 51.45 & 44.75 & 40.83 & 29.17 & 26.35 & \multicolumn{1}{c|}{30.94} & 35.00 & 34.38 & 40.15\\
(\begin{CJK}{UTF8}{gbsn} {文心一言}\end{CJK})  \\
Zhongjing \\
Claude \\
ChatGPT & 39.33 & 35.00 & \multicolumn{1}{c|}{38.79} & 35.33 & 39.33 & 39.77 & 59.40 & 49.18 & \multicolumn{1}{c|}{47.95} & 51.00 & 46.22 & 47.94 & 23.81 & 30.54 & \multicolumn{1}{c|}{31.39} & 40.00 & 37.29 & 40.72\\
GPT-4& 47.33 & 48.18 & \multicolumn{1}{c|}{53.45} & 50.33 & \textbf{53.67} & \textbf{54.19} & \textbf{79.91} & \textbf{72.54} & \multicolumn{1}{c|}{\textbf{70.90}} & \textbf{74.83} & \textbf{73.11} & \textbf{68.35} & 41.07 & 43.71 & \multicolumn{1}{c|}{47.98} & \textbf{58.33} & \textbf{60.00} & \textbf{58.70}\\ \hline
% HuatuoGPT-II & \textbf{71.96} & \textbf{70.90} & \multicolumn{1}{c|}{\textbf{75.00}} & \textbf{73.33} & 70.66 & 70.46 & 79.91 & 77.95 & \multicolumn{1}{c|}{75.00} & 69.93 & 67.44 & 72.24 & \textbf{57.14} & \textbf{54.49} & \multicolumn{1}{c|}{\textbf{57.39}} & 62.43 & - & 52.48\\
HuatuoGPT-II & \textbf{69.16} & \textbf{69.55} & \multicolumn{1}{c|}{\textbf{71.55}} & \textbf{73.67} & 72.33 & 71.40 & 77.78 & 77.05 & \multicolumn{1}{c|}{67.62} & 73.94 & 67.23 & 73.62 & \textbf{57.14} & \textbf{51.50} & \multicolumn{1}{c|}{\textbf{59.19}} & 64.17 & 61.67 & 52.48\\
\hline


\end{tabular}
\caption{The results of Chinese Medical Licensing Examination Results.}
\end{table*}


\begin{table*}[h]
\centering
\begin{tabular}{|c|c|c|c|}
\hline
Models &  \makecell{Step1 \\ (6308)}  & \makecell{Step2\&3 \\ (5148)}  &  \makecell{Average score \\ (11456)} \\ \hline
ChatGLM3-6B  & 33.39 & 32.49 & 32.98 \\
Baichuan2-7B-Chat       & 38.36        & 37.57           & 38.00                 \\ 
Baichuan2-13B-Chat      & 44.90        & 45.80           & 45.30                 \\ 
LLaMA2-7B-Chat & 35.18 & 35.89 & 35.50 \\
LLaMA2-13B-Chat& 40.55 & 39.73 & 40.18 \\
PMC\_LLaMA\_13B \\ \hline
HuatuoGPT-II            & 52.08        & 54.10           & 53.00                 \\ \hline
\end{tabular}
\caption{The results of The United States Medical Licensing Examination (USMLE) from MedQA.}
\end{table*}


\junying{Upcoming addition: The results of The United States Medical Licensing Examination (USMLE)}


\subsection{Evaluation for Instruction Following}
    \junying{ vs. ChatGPT win rate}
    \paragraph{Evaluation with GPT-4}

    \paragraph{Evaluation with Human}


\begin{table*}[ht]
    \centering
    \caption{HuatuoGPT-II Evaluation using GPT-4 Automatic Assessment. Test set and prompts were from HuatuoGPTv1.}
    \label{HuatuoGPTComparison}
    \begin{tabular}{|l|c|c|c|}
    \hline
    \textbf{HuatuoGPT-II Win Rate}             & \textbf{Win} & \textbf{Tie} & \textbf{Fail} \\
    \hline
    HuatuoGPT-II(7B) vs HuatuoGPT-I   & 82  & 13  & 5    \\
    \hline
    HuatuoGPT-II(7B) vs Baichuan2-13B & 64  & 10  & 26   \\
    \hline
    HuatuoGPT-II(7B) vs GPT4          & 65  & 12  & 23   \\
    \hline
    \end{tabular}
    \end{table*}

    \paragraph{Evaluation on Humans} 


    \benyou{design a larger set and use both GPT4 doctors for evaluation}

\subsection{Medical Knowledge Q\&A}
    \paragraph{PubMedQA}
    The results of PubMedQA.
    
    \paragraph{Truthful-QA}
    The results of medical Truthful-QA.

\subsection{Scaling Law of Pretraining and SFT}
    \paragraph{Pretrained Data vs. SFT Data}
    A comparative study between the responses generated by ChatGPT and those produced using pretrained data, when subjected to the same questions, and their impact on post-training models.

\subsection{Insights into Data Integration}
    \paragraph{Training Objectives}
    Assessment of whether adapting test data into SFT-like instruction enhances training results and retains language model capabilities.
    
    \paragraph{Language}
    Testing if converting multilingual data into monolingual data can bolster knowledge transfer.
    
    \paragraph{Style}
    A study on whether maintaining a uniform data style improves learning results.

\subsection{The Imperative of Unified Training}
    \paragraph{Two-Stage vs. Single-Stage Training}
    A comparative analysis of performance outcomes from two-stage training versus single-stage training.
    
    \paragraph{Detrimental Effects of Two-Stage Training on PLM.}
    Experiments were conducted to test the out-of-domain capacities of two-stage transfer versus one-stage transfer. This out-of-domain capacity indicates the remaining ability of the Pretrained Language Model (PLM).


\subsection{Data Analysis}
    \paragraph{Content Deviation in Data Integration}
    Examination of knowledge retention when data undergoes transformations.

    \paragraph{Characteristics of Different Data Sources}
    Investigating the unique features and implications of various data origins.




\benyou{@junying please give detailed plan for analysis, like case study and ablation or any interetsting results that you want to show}



\section{Related Work}

\benyou{@jiangfeng, please finish it as soon as possible}



Recent research has also indicated that such domain adaption using further training~\cite{cheng2023adapting} leads to a drastic drop despite still benefiting fine-tuning evaluation and knowledge probing tests. This inspires us to design a different protocol for domain adaption. 

Also, \cite{gunasekar2023textbooks,li2023textbooks} shows a possibility that 10B well-selected dataset could achieve a comparable performance to much larger model. 

%%
%% The acknowledgments section is defined using the "acks" environment
%% (and NOT an unnumbered section). This ensures the proper
%% identification of the section in the article metadata, and the
%% consistent spelling of the heading.


\section{Conclusion}

The rapid development of large-scale general models like LLama2 a has spurred interest in creating specialized models for vertical domains, aiming to achieve comparable results with medium-sized models through domain-specific knowledge. The conventional two-step adaptation process, involving further pre-training and supervised fine-tuning, faces challenges such as inconsistent data distribution and the necessity for additional hyper-parameters. To overcome these challenges, we introduced a simplified domain adaptation protocol that unifies the two stages into one, transforming diverse data into straightforward input-output pairs. However, this approach encounters the challenge of handling heterogeneous training data in terms of language and genres. By employing a back-translation method, we  transformed varied data into a uniform format in Chinese, showcasing the potential for extension to other languages.  Say something more

\newpage
\begin{acks}
To Robert, for the bagels and explaining CMYK and color spaces.
\end{acks}

%%
%% The next two lines define the bibliography style to be used, and
%% the bibliography file.
\bibliographystyle{ACM-Reference-Format}
\bibliography{sample-base}

%%
%% If your work has an appendix, this is the place to put it.
\appendix

\section{Research Methods}

\subsection{Part One}

Lorem ipsum dolor sit amet, consectetur adipiscing elit. Morbi
malesuada, quam in pulvinar varius, metus nunc fermentum urna, id
sollicitudin purus odio sit amet enim. Aliquam ullamcorper eu ipsum
vel mollis. Curabitur quis dictum nisl. Phasellus vel semper risus, et
lacinia dolor. Integer ultricies commodo sem nec semper.

\subsection{Part Two}

Etiam commodo feugiat nisl pulvinar pellentesque. Etiam auctor sodales
ligula, non varius nibh pulvinar semper. Suspendisse nec lectus non
ipsum convallis congue hendrerit vitae sapien. Donec at laoreet
eros. Vivamus non purus placerat, scelerisque diam eu, cursus
ante. Etiam aliquam tortor auctor efficitur mattis.

\section{Online Resources}

Nam id fermentum dui. Suspendisse sagittis tortor a nulla mollis, in
pulvinar ex pretium. Sed interdum orci quis metus euismod, et sagittis
enim maximus. Vestibulum gravida massa ut felis suscipit
congue. Quisque mattis elit a risus ultrices commodo venenatis eget
dui. Etiam sagittis eleifend elementum.

Nam interdum magna at lectus dignissim, ac dignissim lorem
rhoncus. Maecenas eu arcu ac neque placerat aliquam. Nunc pulvinar
massa et mattis lacinia.


\end{document}
% \end{CJK}
% \end{CJK}
% \endinput
%%
%% End of file `sample-authordraft.tex'.
